% ====================================================================
\chapter{Solving Inequalities}\label{ch:ineq}
% ====================================================================

\begin{chapterintro}
Inequalities are everywhere in economics: a budget must not be exceeded, output cannot be negative, a
project is worth undertaking if its benefits exceed its costs. Solving an inequality means finding
\emph{every} value of the unknown for which it is true. This chapter develops the method in stages ---
linear inequalities, products of factors, quadratics and fractions --- and shows how each step rests on the
properties of the real numbers from \cref{ch:numbers}.
\end{chapterintro}

\begin{prerequisites}
Equivalence and why it matters for solving (\cref{sec:equivalence}); set-builder notation, unions and
intersections (\cref{ch:sets}); the properties O1, A1--A5 and M1--M6 and the sign rules of \cref{prop:signs}
(\cref{ch:numbers}); interval notation (\cref{sec:intervals}).
\end{prerequisites}

\section{What it means to solve an inequality}\label{sec:solve-meaning}

An inequality in an unknown $x$, such as $2x+1>0$, is a proposition depending on $x$. Its
\term{solution set} is the set of all real $x$ for which it is true, $\{x\in\R\mid 2x+1>0\}$. To
\emph{solve} the inequality is to describe this set as simply as possible, usually as an interval or a union of
intervals.

\begin{core}[Solve by a chain of equivalences]
Transform the inequality step by step, making sure that each step is an \emph{equivalence}
(\cref{sec:equivalence}). Then the final, simple inequality has exactly the same solution set as the
original, and its solution set is the answer. The two tools are:
\begin{itemize}
  \item \textbf{adding} the same number to both sides, which is always reversible (A5, applied twice);
  \item \textbf{multiplying} both sides by the same number, which is reversible if the number is
  non-zero, and which keeps the direction if the number is positive and reverses it if the number is
  negative (M6).
\end{itemize}
\end{core}

\begin{note}[How much justification to write]
The slides justify each step by naming the property used, in order to show that the methods rest on the
properties of the real numbers. They also say: ``When solving problems, it suffices just to do correct
calculations.'' \src{L1, slide 22} In your own answers you do not need to cite A5 or M4 at every line; but you
must get the direction right, and when you multiply or divide by something whose sign matters you should say
what that sign is.
\end{note}

\section{Linear inequalities}\label{sec:linear-ineq}

\begin{example}[The first inequality, step by step][ex:lin-ineq]
Find all $x$ that satisfy $2x+1>0$. \src{L1, slides 22--23}
\solution
\emph{Step 1: remove the constant.} Using A5 and A4, add the number $-1$ to both sides; then use A2--A4 to
simplify the left-hand side, and A1, A3 to find $0+(-1)=-1$:
\[
2x+1>0\ \Longrightarrow\ 2x+1+(-1)>0+(-1)\ \Longrightarrow\ 2x>-1.
\]
The reverse implication works as well: add $1$ to both sides of $2x>-1$,
\[
2x>-1\ \Longrightarrow\ 2x+1>-1+1\ \Longrightarrow\ 2x+1>0.
\]
Therefore the two inequalities are equivalent: $2x+1>0\Leftrightarrow 2x>-1$.

\emph{Step 2: remove the coefficient.} By M4 there is an inverse $1/2$ of the non-zero number $2$, and
$1/2>0$ (\cref{prop:signs}(v)). Using M6, multiply both sides of $2x>-1$ by $1/2$; the direction is kept:
\[
2x>-1\ \Longrightarrow\ \tfrac12\cdot 2x>\tfrac12\cdot(-1)\ \Longrightarrow\ x>-\tfrac12.
\]
Similarly the reverse implication holds, multiplying by $2>0$: $x>-\tfrac12\Rightarrow 2x>-1$.

\emph{Conclusion.} All three inequalities are equivalent,
\[
2x+1>0\ \Longleftrightarrow\ 2x>-1\ \Longleftrightarrow\ x>-\tfrac12,
\]
and the solution of $2x+1>0$ is $x\in(-\tfrac12,\infty)$.
\end{example}

The example shows why both directions are checked. If we had shown only $2x+1>0\Rightarrow x>-\tfrac12$, we
would know that every solution lies in $(-\tfrac12,\infty)$, but not that every number in $(-\tfrac12,\infty)$
is a solution.

\subsection{The general linear inequality}

\begin{proposition}[Solving $ax+b>0$][prop:linear]
Let $a,b\in\R$. The solution set of $ax+b>0$ is
\[
\begin{cases}
\left(-\dfrac ba,\ \infty\right) & \text{if } a>0,\\[2ex]
\left(-\infty,\ -\dfrac ba\right) & \text{if } a<0,\\[2ex]
\R\ \text{ if } b>0,\quad\text{and}\quad \emptyset\ \text{ if } b\le0, & \text{if } a=0.
\end{cases}
\]
\end{proposition}
\emph{Proof.} Adding $-b$ to both sides gives the equivalent inequality $ax>-b$. If $a>0$, multiply by
$1/a>0$, which keeps the direction: $x>-b/a$. If $a<0$, multiply by $1/a<0$, which reverses it: $x<-b/a$.
Both steps are reversible (multiply back by $a$). If $a=0$ the inequality reads $0>-b$, i.e.\ $b>0$, which does
not involve $x$ at all: it is true for every $x$ if $b>0$ and for no $x$ if $b\le0$. \hfill$\square$

\begin{example}[Dividing by a negative number]
Solve $5-3x>11$.
\solution
Add $-5$ to both sides: $-3x>6$. Now divide by $-3$, i.e.\ multiply by $-\tfrac13$. Since $-\tfrac13<0$, the
second clause of M6 applies and the direction \emph{reverses}: $x<-2$. The solution is $x\in(-\infty,-2)$.

\emph{Check.} $x=-3$ gives $5+9=14>11$, true; $x=0$ gives $5>11$, false. Both agree with the answer.
\end{example}

\begin{example}[Collecting terms first]
Find all $x$ satisfying $\tfrac12x-3>2x+1$.
\solution
Add $-2x$ and $+3$ to both sides (A5, twice): $\tfrac12x-2x>1+3$, that is, $-\tfrac32x>4$. Multiply by
$-\tfrac23<0$; the direction reverses (M6): $x<-\tfrac83$. The solution is $x\in(-\infty,-\tfrac83)$.
\emph{Check:} $x=-3$ gives $-4.5>-5$, true; $x=0$ gives $-3>1$, false. \src{SG, mock Q2(b)}
\end{example}

\subsection{Inequalities with a parameter}

A \term{parameter} is a letter standing for a fixed but unspecified number, such as a price or a tax rate.
When an inequality involves a parameter, the answer usually depends on it, and the dependence arises exactly
at the step where we multiply or divide by an expression containing it.

\begin{example}[Solving $ax>1$][ex:ax-gt-1]
Let $a\in\R$ be a parameter. Find all $x$ with $ax>1$.
\solution
The step we would like to take is ``divide by $a$''. Whether this is allowed, and what it does to the
direction, depends on the sign of $a$, so we split into three exhaustive cases (O1).
\begin{itemize}
  \item $a>0$: $1/a$ exists (M4) and is positive, so the direction is kept: $x>1/a$, i.e.\ $x\in(1/a,\infty)$.
  \item $a<0$: $1/a$ exists and is negative, so the direction reverses: $x<1/a$, i.e.\ $x\in(-\infty,1/a)$.
  \item $a=0$: there is no $1/0$, so we cannot divide. The inequality reads $0>1$, false for every $x$: the
  solution set is $\emptyset$.
\end{itemize}
\end{example}

\begin{example}[A parameter multiplying a bracket]
Let $a\in\R$. Find all $x$ satisfying $a(x-1)>0$. \src{SG, mock Q2(d)}
\solution
\begin{itemize}
  \item $a>0$: multiply by $1/a>0$: $x-1>0$, so $x\in(1,\infty)$.
  \item $a<0$: multiply by $1/a<0$, reversing: $x-1<0$, so $x\in(-\infty,1)$.
  \item $a=0$: the left side is $0$ for every $x$, and $0>0$ is false, so the solution set is $\emptyset$.
\end{itemize}
\end{example}

\begin{examtip}[``Your answer should depend on the parameter'']
Parameter questions appeared on all three past ECN115 papers summarised in the study guide, usually with a
hint such as ``consider three cases: $a>0$, $a=0$, $a<0$''. The marks are for carrying out each case correctly
and for not forgetting the case $a=0$, which can never be folded into the other two because $1/0$ does not
exist.
\end{examtip}

\section{Products: sign analysis}\label{sec:sign-analysis}

\begin{example}[A product of two linear factors][ex:product-ineq]
Find all $x$ that satisfy $(2x+1)(x-1)>0$. \src{L1, slide 24}
\solution
Neither $2x+1$ nor $x-1$ can be equal to zero: otherwise $(2x+1)(x-1)=0$ (\cref{prop:zero-times}), and $0>0$
is false. So by trichotomy (O1) each factor is either positive or negative, and it suffices to consider four
cases (\cref{prop:signs}(vi)):
\begin{center}
\begin{tabular}{ccc}
\toprule
$2x+1$ & $x-1$ & $(2x+1)(x-1)$\\
\midrule
$>0$ & $>0$ & $>0$\\
$>0$ & $<0$ & $<0$\\
$<0$ & $>0$ & $<0$\\
$<0$ & $<0$ & $>0$\\
\bottomrule
\end{tabular}
\end{center}
Therefore $(2x+1)(x-1)>0$ if and only if
\[
\big[\,2x+1>0\ \text{ and }\ x-1>0\,\big]\quad\text{or}\quad\big[\,2x+1<0\ \text{ and }\ x-1<0\,\big].
\]
Since $2x+1>0$ is equivalent to $x>-\tfrac12$, $x-1>0$ is equivalent to $x>1$, and $x>1\Rightarrow x>-\tfrac12$,
we have $[2x+1>0\text{ and }x-1>0]\Leftrightarrow x>1$. Similarly $[2x+1<0\text{ and }x-1<0]\Leftrightarrow x<-\tfrac12$.
Thus the solution is
\[
x\in(-\infty,-\tfrac12)\cup(1,\infty).
\]
\end{example}

Two small but important points of logic appear in this example. Inside each case the two conditions are joined
by ``and'', so their solution sets are \emph{intersected}; and a pair of conditions such as ``$x>1$ and
$x>-\tfrac12$'' collapses to the stronger one, $x>1$. The cases themselves are joined by ``or'', so their
solution sets are combined by \emph{union}.

\subsection{The sign chart}

The four-case table can be organised on the number line. The \term{critical points} are the values where a
factor is zero, here $x=-\tfrac12$ and $x=1$. They split the line into intervals, and on each interval every
factor keeps one sign. Record the signs in a chart and multiply.

\begin{figure}[ht]
\centering
\begin{tikzpicture}[x=11mm, y=7mm]
  \draw[-{Stealth}] (-3.5,0) -- (4.5,0) node[right] {$x$};
  \foreach \x/\l in {-1/{-\tfrac12}, 2/{1}} { \draw (\x,0.15) -- (\x,-0.15) node[below] {$\l$}; }
  \node[anchor=east, font=\small] at (-3.6,-1.3) {$2x+1$};
  \node[anchor=east, font=\small] at (-3.6,-2.3) {$x-1$};
  \node[anchor=east, font=\small] at (-3.6,-3.3) {$(2x+1)(x-1)$};
  \foreach \y/\a/\b/\c in {-1.3/$-$/$+$/$+$, -2.3/$-$/$-$/$+$, -3.3/$+$/$-$/$+$} {
    \node at (-2.4,\y) {\a}; \node at (0.5,\y) {\b}; \node at (3.3,\y) {\c};
  }
  \foreach \x in {-1,2} \draw[guide] (\x,-0.4) -- (\x,-3.8);
  \node[font=\small] at (-1,-1.3) {$0$}; \node[font=\small] at (2,-2.3) {$0$};
  \node[font=\small] at (-1,-3.3) {$0$}; \node[font=\small] at (2,-3.3) {$0$};
  \draw[accent, line width=2.4pt] (-3.4,-4.4) -- (-1,-4.4); \draw[accent, line width=2.4pt] (2,-4.4) -- (4.4,-4.4);
  \node[circle, draw=accent, fill=white, thick, inner sep=1.5pt] at (-1,-4.4) {};
  \node[circle, draw=accent, fill=white, thick, inner sep=1.5pt] at (2,-4.4) {};
  \node[anchor=east, font=\small] at (-3.6,-4.4) {solution of $>0$};
\end{tikzpicture}
\caption{Sign chart for $(2x+1)(x-1)$. The critical points $-\tfrac12$ and $1$ divide the line into three
intervals; on each, each factor has a constant sign. The product is positive on the two outer intervals, so
the solution of $(2x+1)(x-1)>0$ is $(-\infty,-\tfrac12)\cup(1,\infty)$. Open dots show that the critical points
are excluded, because the inequality is strict.}
\label{fig:sign-chart}
\end{figure}

\begin{external}[Why the signs are constant between critical points]
A linear factor $mx+c$ with $m\ne0$ changes sign only at its root $-c/m$: by \cref{prop:linear} it is positive on
one side of the root and negative on the other. So between consecutive critical points no factor changes sign,
and neither does the product. This is why testing one convenient point in each interval is enough.
\end{external}

\begin{example}[A product that must be negative]
Solve $(x+3)(2x-7)<0$.
\solution
The critical points are $x=-3$ and $x=\tfrac72$. A product is negative exactly when the factors have opposite
signs:
\begin{itemize}
  \item $x+3>0$ and $2x-7<0$: $x>-3$ and $x<\tfrac72$, i.e.\ $-3<x<\tfrac72$;
  \item $x+3<0$ and $2x-7>0$: $x<-3$ and $x>\tfrac72$, which is impossible.
\end{itemize}
The solution is $x\in(-3,\tfrac72)$.
\end{example}

\begin{example}[A non-strict inequality][ex:nonstrict]
Find all $x$ satisfying $(3-x)(x+2)\ge0$. \src{SG, mock Q2(c)}
\solution
Now the product may also be \emph{zero}, which happens at $x=3$ and $x=-2$; both satisfy $\ge0$, so they are
included. Away from these points each factor is positive or negative, and the product is positive when the signs
agree:
\begin{itemize}
  \item $3-x>0$ and $x+2>0$: $x<3$ and $x>-2$, i.e.\ $-2<x<3$;
  \item $3-x<0$ and $x+2<0$: $x>3$ and $x<-2$, impossible.
\end{itemize}
Adding the two points where the product is zero gives $x\in[-2,3]$. Note the \emph{square} brackets.
\end{example}

\begin{warning}[Strict and non-strict inequalities]
For $>0$ or $<0$ the critical points are excluded (round brackets); for $\ge0$ or $\le0$ they are included
(square brackets). Writing $(-2,3)$ instead of $[-2,3]$ in \cref{ex:nonstrict} omits two solutions.
\end{warning}

\section{Quadratic inequalities}\label{sec:quadratic}

A quadratic inequality, such as $3x^2-5x+2>0$, can be solved by rewriting the quadratic as a product of two
linear factors and then using sign analysis.

\begin{example}[Factorising through the roots][ex:quad-ineq]
Find all $x$ that satisfy $3x^2-5x+2>0$. \src{L1, slide 25}
\solution
\emph{Idea.} Represent $3x^2-5x+2$ as a product of two linear terms (if possible),
$3x^2-5x+2=e\cdot(x+f)\cdot(x+g)$ for some real numbers $e,f,g$.

\emph{Finding $e$.} Expanding, $e(x+f)(x+g)=ex^2+e(f+g)x+efg$. Comparing the coefficients of $x^2$ gives $e=3$.
Comparing the others gives only $f+g=-\tfrac53$ and $fg=\tfrac23$, which do not immediately give $f$ and $g$.

\emph{Observation.} Substitute $x=-f$: then $e(x+f)(x+g)=e\cdot0\cdot(g-f)=0$. Thus $x=-f$ is a solution
(\emph{root}) of the equation $3x^2-5x+2=0$, and similarly $x=-g$ is a root. So $f$ and $g$ can be found by
solving the equation.

\emph{Solving the equation.} The discriminant is $D=(-5)^2-4\cdot3\cdot2=25-24=1$, so the roots are
\[
x_{1,2}=\frac{-(-5)\pm\sqrt1}{2\cdot 3}=\frac{5\pm1}{6},\qquad x_1=\tfrac23,\quad x_2=1.
\]
Hence $3x^2-5x+2=3\,(x-\tfrac23)(x-1)$.

\emph{Sign analysis.} Solving $3(x-\tfrac23)(x-1)>0$ as in \cref{ex:product-ineq} (the factor $3$ is positive and
does not affect the sign), the inequality $3x^2-5x+2>0$ is equivalent to
\[
x\in(-\infty,\tfrac23)\cup(1,\infty).
\]
\end{example}

\subsection{The quadratic formula}

\begin{keyformula}[Roots of a quadratic][kf:quadratic]{x_{1,2}=\frac{-b\pm\sqrt{D}}{2a},\qquad D=b^2-4ac}
\fsymbols $a,b,c$ are the real coefficients of $ax^2+bx+c$, with $a\ne0$; $D$ is the \emph{discriminant};
$x_1,x_2$ are the roots, the values of $x$ where $ax^2+bx+c=0$; the symbol $\pm$ means that one root uses $+$
and the other $-$.
\fmeaning The formula gives every real solution of $ax^2+bx+c=0$. The discriminant tells you how many there are:
two distinct real roots if $D>0$, one repeated root $x=-b/(2a)$ if $D=0$, and no real roots if $D<0$ (the square
root of a negative number is not a real number).
\forigin Completing the square: $ax^2+bx+c=a\big[(x+\tfrac b{2a})^2-\tfrac{D}{4a^2}\big]$ (\cref{kf:complete-square}).
Setting this to zero gives $(x+\tfrac b{2a})^2=\tfrac D{4a^2}$, and taking square roots (possible only when $D\ge0$)
gives the formula.
\fuse Whenever a quadratic must be factorised and the factors are not obvious; in particular, as the first step in
solving any quadratic inequality.
\fexample For $3x^2-5x+2$: $a=3$, $b=-5$, $c=2$. Then $D=(-5)^2-4(3)(2)=1$, $\sqrt D=1$, and
$x=\frac{5\pm1}{6}$, giving $\tfrac23$ and $1$. Check: $3(\tfrac23)^2-5(\tfrac23)+2=\tfrac43-\tfrac{10}3+2=0$.
\fmistakes Writing $-b$ incorrectly when $b$ is negative (here $-b=+5$); dividing only $\sqrt D$ by $2a$ instead
of the whole numerator; forgetting that $D<0$ means there are no real roots, not that the formula ``fails''.
\frearrange Once the roots are known, the quadratic factorises as
$ax^2+bx+c=a(x-x_1)(x-x_2)$. Comparing coefficients gives $x_1+x_2=-b/a$ and $x_1x_2=c/a$, a quick check on
the roots.
\fchanges Increasing $c$ (with $a>0$) lowers $D=b^2-4ac$: the two roots move towards each other, merge when
$D=0$, and disappear when $D<0$. Geometrically the parabola is lifted upwards until it no longer meets the
$x$-axis.
\fexam Usually as an unmarked first step (``factorise'', ``find the roots''). A parameter in $c$ leads to a question
such as ``for which values of $c$ does the equation have two distinct roots?'', answered by solving $D>0$.
\end{keyformula}

\subsection{The sign of a quadratic}

Once $ax^2+bx+c=a(x-x_1)(x-x_2)$ with $x_1<x_2$, sign analysis gives a rule worth remembering.

\begin{proposition}[Sign of a quadratic with two roots][prop:quad-sign]
Let $a>0$ and $x_1<x_2$. Then $a(x-x_1)(x-x_2)$ is positive for $x<x_1$ and for $x>x_2$, zero at $x_1$ and $x_2$,
and negative for $x_1<x<x_2$. If $a<0$ all the signs are reversed.
\end{proposition}
\emph{Proof.} For $x<x_1$ both factors are negative; for $x>x_2$ both are positive; between the roots
$x-x_1>0$ and $x-x_2<0$. Apply \cref{prop:signs}(vi), then multiply by $a$ using M6. \hfill$\square$

\begin{figure}[ht]
\centering
\begin{tikzpicture}
\begin{axis}[stat, width=9cm, height=6.2cm, xmin=-0.6, xmax=2.2, ymin=-0.6, ymax=3.2,
  xlabel={$x$}, ylabel={$y$}, xtick={0,0.6667,1,2}, xticklabels={$0$,$\tfrac23$,$1$,$2$},
  ytick={0,1,2,3}, axis x line=middle, axis y line=middle, x label style={anchor=west}, y label style={anchor=south}]
  \addplot[curve, accent, domain=-0.45:2.05, samples=100] {3*x^2-5*x+2};
  \addplot[line width=3pt, corecol, opacity=0.8, domain=-0.6:0.6667] {0};
  \addplot[line width=3pt, corecol, opacity=0.8, domain=1:2.2] {0};
  \node[circle, draw=corecol, fill=white, thick, inner sep=1.5pt] at (axis cs:0.6667,0) {};
  \node[circle, draw=corecol, fill=white, thick, inner sep=1.5pt] at (axis cs:1,0) {};
  \node[lab, accent, anchor=west] at (axis cs:0.15,2.6) {$y=3x^2-5x+2$};
\end{axis}
\end{tikzpicture}
\caption{The parabola $y=3x^2-5x+2$ opens upwards because the leading coefficient $3$ is positive. It crosses the
$x$-axis at the roots $\tfrac23$ and $1$, lies below the axis between them and above it outside them. The thick
segments show the solution set of $3x^2-5x+2>0$.}
\label{fig:parabola}
\end{figure}

\begin{warning}[Answering between the roots]
For a quadratic with positive leading coefficient, ``$>0$'' is solved \emph{outside} the roots and ``$<0$''
\emph{between} them. A quick test of one point (for example $x=0$, if it is not a root) catches the error: in
\cref{ex:quad-ineq}, $x=0$ gives $2>0$, and $0$ is indeed outside $[\tfrac23,1]$.
\end{warning}

\subsection{When the discriminant is zero or negative}

The lecture example has $D>0$. The other two cases are handled by completing the square.

\begin{keyformula}[Completing the square][kf:complete-square]{ax^2+bx+c=a\left(x+\frac{b}{2a}\right)^2-\frac{D}{4a},\qquad D=b^2-4ac}
\fsymbols $a\ne0$, $b$, $c$ real coefficients; $D$ the discriminant.
\fmeaning Every quadratic is a multiple of a perfect square plus a constant. Since a square is never negative, the
sign of the quadratic can be read off directly.
\forigin Expand the right-hand side:
$a\big(x^2+\tfrac bax+\tfrac{b^2}{4a^2}\big)-\tfrac{b^2-4ac}{4a}=ax^2+bx+\tfrac{b^2}{4a}-\tfrac{b^2}{4a}+c=ax^2+bx+c$.
\fuse When $D=0$ or $D<0$, or when you need the minimum (or maximum) value of a quadratic.
\fexample $x^2-5x+\tfrac{25}{4}=(x-\tfrac52)^2$ ($D=0$), and $x^2-5x+7=(x-\tfrac52)^2+\tfrac34$ ($D=-3$).
\fmistakes Forgetting to factor out $a$ before completing the square; sign errors in $\tfrac b{2a}$.
\fchanges If $a>0$ the quadratic is smallest at $x=-\tfrac b{2a}$, where it equals $-\tfrac D{4a}$. If $D<0$ this
minimum is positive, so the quadratic is positive everywhere; if $D=0$ the minimum is $0$, attained only at
$x=-\tfrac b{2a}$.
\fexam ``Show that the expression is positive for all $x$'', with the hint ``write it as a square plus a constant''.
\end{keyformula}

\begin{proposition}[Sign of a quadratic when $D\le0$][prop:quad-D0]
Let $a>0$. If $D=0$, then $ax^2+bx+c=a(x-x_0)^2$ with $x_0=-\tfrac b{2a}$, which is positive for every
$x\ne x_0$ and zero at $x_0$. If $D<0$, then $ax^2+bx+c>0$ for every real $x$.
\end{proposition}
\emph{Proof.} By \cref{kf:complete-square}, $ax^2+bx+c=a(x-x_0)^2-\tfrac D{4a}$. A square is $\ge0$, and is $>0$
unless what is squared is zero (\cref{prop:signs}(iv)). If $D=0$ the constant vanishes; if $D<0$ the constant
$-\tfrac D{4a}$ is positive, so the sum is positive. \hfill$\square$

\begin{example}[A quadratic with a parameter][ex:quad-param]
Let $c\in\R$ and consider $x^2-5x+c$. \src{SG, mock Q3}
\begin{enumerate}
  \item For which $c$ does $x^2-5x+c=0$ have two distinct real roots?
  \item Solve $x^2-5x+6>0$.
  \item Solve $x^2-5x+\tfrac{25}4>0$.
  \item Show that if $c>\tfrac{25}{4}$ then $x^2-5x+c>0$ for every $x\in\R$.
\end{enumerate}
\solution
(a) $D=(-5)^2-4\cdot1\cdot c=25-4c$. Two distinct roots exactly when $D>0$: $25-4c>0\Leftrightarrow c<\tfrac{25}4$
(dividing by $4>0$ keeps the direction). So $c\in(-\infty,\tfrac{25}4)$.

(b) $D=1$, roots $\frac{5\pm1}2=2,3$, so $x^2-5x+6=(x-2)(x-3)$. The leading coefficient is positive, so the
expression is positive outside the roots: $x\in(-\infty,2)\cup(3,\infty)$.

(c) $D=0$ and $x^2-5x+\tfrac{25}4=(x-\tfrac52)^2$, which is positive except at $x=\tfrac52$. So
$x\in(-\infty,\tfrac52)\cup(\tfrac52,\infty)=\R\setminus\{\tfrac52\}$. In (b) a whole interval was excluded; here
the two roots have merged and only one point is excluded.

(d) Completing the square, $x^2-5x+c=(x-\tfrac52)^2+(c-\tfrac{25}4)$. The square is $\ge0$ and, when
$c>\tfrac{25}4$, the constant is $>0$; so the sum is positive for every $x$. (This is the case $D<0$.)
\end{example}

\begin{example}[A non-strict quadratic inequality]
Solve $2x^2-7x+3\ge0$.
\solution
$D=49-24=25$, roots $\frac{7\pm5}{4}$, i.e.\ $\tfrac12$ and $3$. So $2x^2-7x+3=2(x-\tfrac12)(x-3)$. With leading
coefficient $2>0$, the expression is $\ge0$ outside the roots and at the roots themselves:
$x\in(-\infty,\tfrac12]\cup[3,\infty)$.
\end{example}

\section{Inequalities involving fractions}\label{sec:fraction-ineq}

\begin{example}[An incorrect argument][ex:one-over-x]
A student solves $\dfrac1x>2$ as follows: ``Multiply both sides by $x$ to obtain $1>2x$. Dividing by $2$ gives
$x<\tfrac12$. Therefore the solution is $x\in(-\infty,\tfrac12)$.'' Explain the error and find the correct
solution. \src{SG, mock Q6}
\solution
\emph{The error.} Multiplying an inequality by $x$ is governed by M6, which keeps the direction only if $x>0$ and
reverses it if $x<0$. The sign of $x$ is unknown, so the step is not justified. Moreover $x=0$ is not allowed at
all, since $1/0$ is undefined.

\emph{A value wrongly included.} $x=-1$ lies in the student's answer, but $1/(-1)=-1>2$ is false.

\emph{Correct solution.} Split by the sign of $x$.
\begin{itemize}
  \item $x>0$: multiplying by $x>0$ keeps the direction, $1>2x$, so $x<\tfrac12$. With $x>0$: $0<x<\tfrac12$.
  \item $x<0$: then $1/x<0<2$ (\cref{prop:signs}(v)), so the inequality fails for every negative $x$.
  \item $x=0$: excluded, as $1/x$ is undefined.
\end{itemize}
The solution set is $(0,\tfrac12)$.
\end{example}

The safest general method avoids multiplying by an expression of unknown sign altogether.

\begin{core}[Method for inequalities with fractions]
\begin{enumerate}[label=\arabic*.]
  \item Move everything to one side, so that the other side is $0$.
  \item Combine into a single fraction $\dfrac{N(x)}{M(x)}$ and factorise numerator and denominator.
  \item Apply sign analysis to the factors of $N$ and $M$: a quotient has the same sign as the corresponding
  product, because $N/M=N\cdot M\cdot(1/M^2)$ and $1/M^2>0$.
  \item Exclude the zeros of $M$ (the fraction is undefined there); include the zeros of $N$ only for
  $\ge$ or $\le$.
\end{enumerate}
\end{core}

\begin{example}[Using the method]
Solve $\dfrac1x>2$ and $\dfrac{x+1}{x-2}\ge0$.
\solution
For the first, $\frac1x-2>0\Leftrightarrow\frac{1-2x}{x}>0$. Critical points $x=0$ and $x=\tfrac12$. Signs:
for $x<0$, numerator $+$, denominator $-$, quotient $-$; for $0<x<\tfrac12$, both $+$, quotient $+$; for
$x>\tfrac12$, numerator $-$, denominator $+$, quotient $-$. Solution $(0,\tfrac12)$, as before.

For the second, critical points $x=-1$ (numerator zero) and $x=2$ (denominator zero). For $x<-1$ both factors
are negative, quotient positive; for $-1<x<2$ the quotient is negative; for $x>2$ it is positive. Include $x=-1$
(where the quotient is $0$) but not $x=2$ (undefined): $x\in(-\infty,-1]\cup(2,\infty)$.
\end{example}

\section{Inequalities in economics}\label{sec:ineq-econ}

\begin{example}[Break-even output][ex:break-even]
A firm sells its product at price $p=12$ per unit. Its total cost of producing $q$ units is $C(q)=200+7q$
(a fixed cost of $200$ and a cost of $7$ per unit). For which outputs $q\ge0$ does the firm make a strictly
positive profit?
\solution
Profit is revenue minus cost: $\pi(q)=12q-(200+7q)=5q-200$. We need $5q-200>0\Leftrightarrow 5q>200\Leftrightarrow q>40$
(multiplying by $\tfrac15>0$). The firm makes a profit exactly when $q\in(40,\infty)$. The value $q=40$ is the
\emph{break-even output}.
\end{example}

\begin{example}[A quadratic profit function]
Suppose instead that profit is $\pi(q)=-q^2+10q-16$. For which $q$ is profit positive?
\solution
$-q^2+10q-16>0\Leftrightarrow q^2-10q+16<0$ (multiplying by $-1$ reverses the direction). The roots of
$q^2-10q+16$ are $\frac{10\pm\sqrt{100-64}}2=\frac{10\pm6}2$, i.e.\ $2$ and $8$, so $q^2-10q+16=(q-2)(q-8)$, which is
negative between the roots. Profit is positive exactly for $q\in(2,8)$: producing too little fails to cover
fixed costs, and producing too much drives costs up faster than revenue.
\end{example}

% --------------------------------------------------------------------
\section{Chapter review}
% --------------------------------------------------------------------
\subsection*{Summary}
Solving an inequality means describing its solution set exactly, which requires a chain of equivalences.
Adding a number to both sides is always reversible; multiplying by a positive number keeps the direction,
by a negative number reverses it, and by zero destroys it (M6). A linear inequality $ax+b>0$ is solved by one
addition and one multiplication, with three cases when $a$ is a parameter. A product is positive when its
factors share a sign and negative when they have opposite signs; organising this on the number line around the
critical points gives the sign chart. A quadratic is factorised through its roots, found with the discriminant;
with positive leading coefficient it is positive outside its roots and negative between them, and completing the
square handles the cases $D=0$ and $D<0$. Fractions are handled by moving everything to one side and analysing
the signs of numerator and denominator, never by multiplying by an expression of unknown sign.

\subsection*{Key definitions}
\begin{itemize}
  \item \textbf{Solution set} of an inequality: the set of all values of the unknown for which it is true.
  \item \textbf{Critical points}: the values at which a factor (of a numerator or denominator) is zero.
  \item \textbf{Discriminant} of $ax^2+bx+c$: $D=b^2-4ac$.
\end{itemize}

\subsection*{Key equations}
\begin{gather*}
D=b^2-4ac,\qquad x_{1,2}=\frac{-b\pm\sqrt D}{2a},\\
ax^2+bx+c=a(x-x_1)(x-x_2)\ \ (D\ge0),\qquad
ax^2+bx+c=a\Big(x+\frac b{2a}\Big)^2-\frac D{4a}.
\end{gather*}

\subsection*{Connections}
Every step uses the properties of \cref{ch:numbers}. Answers are sets, written with the interval and union
notation of \cref{ch:sets}. Absolute-value inequalities (\cref{ch:abs}) reduce to the methods here, and the
definition of a limit (\cref{ch:limits}) requires solving inequalities such as $2^{-n}<\varepsilon$ for $n$.
In Topic~4, the sign of a derivative written as a fraction $P(x)/Q(x)$ is found by exactly this sign analysis.

\subsection*{Common mistakes}
\begin{itemize}
  \item Not reversing the inequality when multiplying or dividing by a negative number.
  \item Dividing by a parameter or variable without considering its sign, or forgetting the zero case.
  \item Multiplying both sides by $x$ in an inequality with $1/x$.
  \item Using $\cap$ to combine the cases of a sign analysis (cases are combined with $\cup$).
  \item Solving ``$>0$'' for a quadratic with positive leading coefficient by taking the interval between the roots.
  \item Wrong brackets: round for strict inequalities and for excluded points, square for included endpoints.
\end{itemize}

\subsection*{Exam checklist}
\begin{checklist}
  \item Solve a linear inequality, stating where the direction changes and why.
  \item Solve $ax>b$ or $a(x-k)>0$ for all values of a parameter $a$.
  \item Carry out a sign analysis for a product of linear factors, including non-strict cases.
  \item Find roots with the discriminant and factorise a quadratic.
  \item Solve quadratic inequalities for $D>0$, $D=0$ and $D<0$.
  \item Solve inequalities involving fractions without multiplying by an expression of unknown sign.
\end{checklist}

% --------------------------------------------------------------------
\section{Practice questions}
% --------------------------------------------------------------------
\pq[basic] Solve: (a) $3x-4\le 8$; (b) $7-2x>1$; (c) $\tfrac{x}{3}+1\ge \tfrac{x}{2}$.

\pq[basic] Solve $(x-4)(x+1)>0$ and $(x-4)(x+1)\le0$.

\pq[mathematical] Solve $x^2-x-6<0$ and $4x^2+4x+1>0$.

\pq[mathematical] Show that $x^2+2x+5>0$ for every real $x$.

\pq[conceptual] Explain why, in a sign analysis, the cases are combined with $\cup$ while the conditions within a
case are combined with $\cap$.

\pq[mathematical] Solve $\dfrac{2x-1}{x+3}<1$.

\pq[application] A consumer has $\pounds 60$ to spend on two goods priced $\pounds 4$ and $\pounds 6$ per unit. She
wants to buy $x$ units of the first good and exactly $6$ units of the second. For which $x\ge0$ can she afford
this?

\pq[multi-step] Let $b\in\R$ be a parameter. Find all $x$ with $bx+3>x$. Your answer should depend on $b$.

\pq[multi-step] For which values of the parameter $k$ is $x^2+kx+4>0$ for every real $x$?

\pq[exam-style] Let $a\in\R$ be a parameter.
(a) Find all $x$ satisfying $(x-a)(x-1)<0$. Your answer should depend on $a$. Hint: consider the cases $a>1$,
$a=1$ and $a<1$.
(b) A student claims that $\dfrac{x-1}{x}\ge1$ is equivalent to $x-1\ge x$ and hence has no solutions. Explain
why the argument is not correct and find the solution set.

% --------------------------------------------------------------------
\section{Solutions to practice questions}
% --------------------------------------------------------------------
\pqsol{4.1}
(a) $3x\le12$, so $x\le4$: $x\in(-\infty,4]$.
(b) $-2x>-6$; dividing by $-2<0$ reverses: $x<3$, so $x\in(-\infty,3)$.
(c) Multiply by $6>0$: $2x+6\ge3x$, so $6\ge x$: $x\in(-\infty,6]$.

\pqsol{4.2}
Critical points $-1$ and $4$. The product is positive when both factors have the same sign: $x<-1$ or $x>4$. So
$(x-4)(x+1)>0\Leftrightarrow x\in(-\infty,-1)\cup(4,\infty)$, and $(x-4)(x+1)\le0\Leftrightarrow x\in[-1,4]$
(between the roots, roots included).

\pqsol{4.3}
$x^2-x-6=(x-3)(x+2)$, negative between the roots: $x\in(-2,3)$.
$4x^2+4x+1=(2x+1)^2$ ($D=0$), positive except at $x=-\tfrac12$: $x\in\R\setminus\{-\tfrac12\}$.

\pqsol{4.4}
$x^2+2x+5=(x+1)^2+4$. A square is $\ge0$, so the sum is $\ge4>0$ for every $x$. (Equivalently $D=4-20<0$ with a
positive leading coefficient.)

\pqsol{4.5}
Within a case, \emph{both} conditions must hold (for instance ``$2x+1>0$ \emph{and} $x-1>0$''), and the set of $x$
satisfying two conditions simultaneously is the intersection of their solution sets. A number solves the
original inequality if it falls into \emph{some} case (``case 1 \emph{or} case 4''), and the set of $x$ satisfying
at least one of several conditions is the union of their solution sets. Intersection corresponds to ``and'' and
union to ``or'' (\cref{sec:intersection-union}).

\pqsol{4.6}
$\frac{2x-1}{x+3}-1<0\Leftrightarrow\frac{2x-1-(x+3)}{x+3}<0\Leftrightarrow\frac{x-4}{x+3}<0$. Critical points
$-3$ and $4$; the quotient is negative when the signs differ, which happens for $-3<x<4$. Solution $(-3,4)$.

\pqsol{4.7}
Cost $4x+6\cdot6=4x+36\le60\Leftrightarrow4x\le24\Leftrightarrow x\le6$. With $x\ge0$: $x\in[0,6]$.

\pqsol{4.8}
$bx+3>x\Leftrightarrow(b-1)x>-3$.
If $b>1$: $x>-\frac3{b-1}$, i.e.\ $x\in\big(-\frac3{b-1},\infty\big)$.
If $b<1$: dividing by $b-1<0$ reverses, $x<-\frac3{b-1}=\frac3{1-b}$, i.e.\ $x\in\big(-\infty,\frac3{1-b}\big)$.
If $b=1$: $0>-3$, true for all $x$: the solution set is $\R$.

\pqsol{4.9}
With leading coefficient $1>0$, $x^2+kx+4>0$ for all $x$ exactly when there are no real roots, $D<0$ (if $D=0$
the quadratic is zero at one point; if $D>0$ it is negative between the roots). $D=k^2-16<0\Leftrightarrow
(k-4)(k+4)<0\Leftrightarrow k\in(-4,4)$.

\pqsol{4.10}
(a) The critical points are $a$ and $1$; the product is negative strictly between them.
If $a>1$: $x\in(1,a)$. If $a<1$: $x\in(a,1)$. If $a=1$: $(x-1)^2<0$ has no solution, so the solution set is
$\emptyset$.
(b) Multiplying by $x$ is valid only for $x>0$; for $x<0$ the direction reverses, and $x=0$ is excluded. Correct
method: $\frac{x-1}{x}-1\ge0\Leftrightarrow\frac{-1}{x}\ge0$. The numerator is negative, so the quotient is
$\ge0$ exactly when $x<0$ (and it is never $0$). Solution set $(-\infty,0)$. For example $x=-1$ gives
$\frac{-2}{-1}=2\ge1$, a solution the student's argument missed.
